\DocumentMetadata{
  pdfversion=2.0,
  pdfstandard=ua-2,
  lang=en-US,
  tagging=on,
  tagging-setup={math/setup=mathml-SE,tabsorder=structure}
}
\documentclass[11pt,letterpaper]{article}
\usepackage[margin=0.68in,includefoot]{geometry}
\usepackage{newcomputermodern}
\usepackage{amsmath,amscd,mathtools}
\usepackage{tikz,adjustbox}
\providecommand{\square}{\mdlgwhtsquare}
\usepackage[hidelinks]{hyperref}
\hypersetup{
  pdftitle={Algebra PhD Exam, September 2003},
  pdfauthor={University of Florida Department of Mathematics},
  pdfsubject={Graduate mathematics examination},
  pdfdisplaydoctitle=true,
  bookmarksopen=true
}
\setcounter{secnumdepth}{0}
\setlength{\parindent}{0pt}
\setlength{\parskip}{0.28em}
\setlength{\emergencystretch}{3em}
\setlength{\leftmargini}{2.15em}
\setlength{\leftmarginii}{2.65em}
\setlength{\leftmarginiii}{3em}
\pagestyle{empty}
\raggedbottom
\allowdisplaybreaks
\NewTaggingSocketPlug{block/list/label}{examproblem}{%
  \tagstructbegin{tag=Lbl}\tagstructend%
  \tagstructbegin{tag=\LBody}%
  \tagstructbegin{tag=\ProblemHeadingTag,title-o={\ProblemHeadingText}}%
  \tagmcbegin{tag=\ProblemHeadingTag,actualtext={\ProblemHeadingText}}%
  #2%
  \tagmcend%
  \tagstructend%
}
\newenvironment{examproblems}[2]{%
  \def\ProblemHeadingTag{#2}%
  \AssignTaggingSocketPlug{block/list/label}{examproblem}%
  \begin{enumerate}%
  \setcounter{enumi}{#1}%
  \renewcommand{\labelenumi}{\textbf{\theenumi.}}%
  \setlength{\topsep}{0.35em}%
  \setlength{\partopsep}{0pt}%
  \setlength{\itemsep}{0.48em}%
  \setlength{\parsep}{0.18em}%
}{\end{enumerate}}
\newenvironment{examsubparts}{%
  \AssignTaggingSocketPlug{block/list/label}{default}%
  \begin{enumerate}%
  \setlength{\topsep}{0.18em}%
  \setlength{\partopsep}{0pt}%
  \setlength{\itemsep}{0.16em}%
  \setlength{\parsep}{0.1em}%
}{\end{enumerate}}
\newenvironment{examinstructions}{%
  \par\begingroup\itshape
}{%
  \par\endgroup\vspace{0.18em}
}
\makeatletter
\renewcommand\subsection{\@startsection{subsection}{2}{0pt}%
  {-1.25ex plus -.25ex minus -.1ex}%
  {0.55ex plus .1ex}%
  {\normalfont\large\bfseries}}
\renewcommand{\@maketitle}{%
  \newpage\null\vskip -1em%
  {\footnotesize\bfseries
    \href{https://www.ufl.edu/}{University of Florida}, \href{https://math.ufl.edu/}{Department of Mathematics}\par}%
  \vskip -0.5em%
  \tagstructbegin{tag=H1,title={Algebra PhD Exam, September 2003}}%
  \tagpdfparaOff%
  \tagmcbegin{tag=H1}%
    {\LARGE\bfseries\href{https://gma.math.ufl.edu/exams/algebra-phd/}{Algebra PhD Exam}, September 2003\par}%
  \tagmcend%
  \tagpdfparaOn%
  \tagstructend%
  \vskip 0.25em%
  \tagmcbegin{artifact}%
    \hrule height 1pt%
  \tagmcend%
  \vskip 1em%
}
\makeatother
\title{Algebra PhD Exam, September 2003}
\author{}
\date{}
\begin{document}
\maketitle
\thispagestyle{empty}
\begin{examinstructions}
Do seven of the following eleven problems. Do not turn in more than seven problems. You must show your work. Answers with no work and/or no explanations will receive no credit. State clearly any theorem you use in your proofs. In the problems, \(\mathbb{Z}\), resp. \(\mathbb{Q}\), \(\mathbb{C}\), is the set of all integers, resp. of all rational numbers, of all complex numbers.
\end{examinstructions}
\begin{examproblems}{0}{H2}
\def\ProblemHeadingText{Problem 1.}
\item
Consider the following statement \(A(P)\): “If a normal subgroup~\(H\) of a group~\(G\) and the quotient \(G/H\) both have property~\(P\), then so does~\(G\).” Prove or disprove \(A(P)\), where
\begin{examsubparts}
\item[\textnormal{a)}]
\(P\) is “being solvable”;
\item[\textnormal{b)}]
\(P\) is “being nilpotent”.
\end{examsubparts}
\def\ProblemHeadingText{Problem 2.}
\item
Inside the symmetric group \(S_7\), consider the subgroup~\(G\) generated by the permutations \((1234567)\) and \((235)(476)\). Let
\[
H=\langle a,b\mid a^7=b^3=1,\ bab^{-1}=a^2\rangle.
\]
 Show that \(G\simeq H\).
\def\ProblemHeadingText{Problem 3.}
\item
Let~\(G\) be a finite group with exactly one maximal subgroup. Prove that~\(G\) is cyclic of prime power order.
\def\ProblemHeadingText{Problem 4.}
\item
Let~\(F\) be a finite field of finite cardinality~\(q\) and~\(n\) be any natural number. Show that there is at least one irreducible polynomial \(f\in F[x]\) of degree~\(n\).
\def\ProblemHeadingText{Problem 5.}
\item
Let~\(A\) be a commutative ring with identity. Suppose that~\(M\) and~\(N\) are free~\(A\)-modules with~\(m\) and~\(n\) generators, respectively. Prove that if \(M\simeq N\), then \(m=n\).
\def\ProblemHeadingText{Problem 6.}
\item
This problem has two parts.
\begin{examsubparts}
\item[\textnormal{a)}]
Give an example of a projective~\(\mathbb{Z}\)-module and an example of an injective~\(\mathbb{Z}\)-module.
\item[\textnormal{b)}]
Let~\(A\) be a finite abelian group considered as a~\(\mathbb{Z}\)-module and let~\(I\) be an injective~\(\mathbb{Z}\)-module. Compute \(A\otimes_{\mathbb{Z}} I\).
\end{examsubparts}
\def\ProblemHeadingText{Problem 7.}
\item
Let \(\alpha\in\mathbb{C}\) be such that \([\mathbb{Q}(\alpha):\mathbb{Q}]=2003\). Let \(E/\mathbb{Q}\) be a normal closure of \(\mathbb{Q}(\alpha)/\mathbb{Q}\). Prove that \([E:\mathbb{Q}]\) divides \(2003!\).
\def\ProblemHeadingText{Problem 8.}
\item
Let \(\xi\in\mathbb{C}\) be a primitive \(77^{\mathrm{th}}\) root of unity in~\(\mathbb{C}\), and let \(F=\mathbb{Q}(\xi)\).
\begin{examsubparts}
\item[\textnormal{a)}]
Briefly explain why \(F/\mathbb{Q}\) is a Galois extension, and describe the structure of the Galois group \(\operatorname{Aut}_{\mathbb{Q}}(F)\).
\item[\textnormal{b)}]
Find the number of subfields of~\(F\) that are quadratic extensions of~\(\mathbb{Q}\).
\end{examsubparts}
\def\ProblemHeadingText{Problem 9.}
\item
State and prove the Hilbert Basis Theorem.
\def\ProblemHeadingText{Problem 10.}
\item
Let~\(D\) be a Dedekind domain and~\(F\) be the group of fractional ideals of~\(D\). Prove that any nontrivial element in~\(F\) has infinite order.
\def\ProblemHeadingText{Problem 11.}
\item
Let~\(R\) be a commutative ring with identity such that every element \(x\in R\) satisfies \(x^{n(x)}=x\) for some integer \(n(x)\geq 1\) depending on~\(x\). Prove that the Jacobson radical of~\(R\) is~\(0\).
\end{examproblems}
\end{document}
